\begin{thebibliography}{99}

\bibitem{ref1}
Perera, A. T. D., Javanroodi, K., et al. (2023). Challenges resulting from urban density and climate change for the EU energy transition. \textit{Nature Energy}, 8(4), 397--412. \url{https://doi.org/10.1038/s41560-023-01232-9}

\bibitem{ref2}
van Heerden, R., Edelenbosch, O. Y., et al. (2025). Demand-side strategies enable rapid and deep cuts in buildings and transport emissions to 2050. \textit{Nature Energy}, 10(3), 380--394. \url{https://doi.org/10.1038/s41560-025-01703-1}

\bibitem{ref6}
Javanroodi, K., Perera, A. T. D., et al. (2023). Designing climate resilient energy systems in complex urban areas considering urban morphology: A technical review. \textit{Advances in Applied Energy}, 12. \url{https://doi.org/10.1016/j.adapen.2023.100155}

\bibitem{ref7}
Perera, A. T. D., Javanroodi, K., \& Nik, V. M. (2021). Climate resilient interconnected infrastructure: Co-optimization of energy systems and urban morphology. \textit{Applied Energy}, 285. \url{https://doi.org/10.1016/j.apenergy.2020.116430}

\bibitem{ref3}
Zhou, B., \& Hu, M. (2025). Impact of architectural design elements on the energy efficiency and thermal comfort in neighborhood building ensembles. \textit{Journal of Building Engineering}, 100. \url{https://doi.org/10.1016/j.jobe.2024.111708}

\bibitem{ref4}
Du, L., Wang, H., et al. (2025). Impact of block form on building energy consumption, urban microclimate and solar potential: A case study of Wuhan, China. \textit{Energy and Buildings}, 328. \url{https://doi.org/10.1016/j.enbuild.2024.115224}

\bibitem{ref5}
Liu, B., Liu, Y., et al. (2024). Urban morphology indicators and solar radiation acquisition: 2011--2022 review. \textit{Renewable and Sustainable Energy Reviews}, 199. \url{https://doi.org/10.1016/j.rser.2024.114548}

\bibitem{ref41}
Wang, H., Wu, Z., Gu, W., et al. (2025). A GIS-informed deep learning framework to assess the effect of urban morphology on energy demand. \textit{Energy and Buildings}, 319, 116535. \url{https://doi.org/10.1016/j.enbuild.2025.116535}

\bibitem{wang2026unlocking}
Wang, H., Wu, Z., Gu, W., et al. (2026). Unlocking the potential of urban spatial structure to shape building energy demand at a large scale. \textit{Sustainable Cities and Society}, 116, 107381. \url{https://doi.org/10.1016/j.scs.2026.107381}

\bibitem{wang2026climateResilient}
Wang, H., Wu, Z., Gu, W., et al. (2026). Towards climate-resilient cities: Assessing and planning urban energy demand under extreme climate events. \textit{Energy and Buildings}, 337, 117436. \url{https://doi.org/10.1016/j.enbuild.2026.117436}

\bibitem{ref8}
Shi, Z., Fonseca, J. A., \& Schlueter, A. (2017). A review of simulation-based urban form generation and optimization for energy-driven urban design. \textit{Building and Environment}, 121, 119--129. \url{https://doi.org/10.1016/j.buildenv.2017.05.006}

\bibitem{ref18}
Pan, Y., Zhu, M., et al. (2023). Building energy simulation and its application for building performance optimization: A review of methods, tools, and case studies. \textit{Advances in Applied Energy}, 10. \url{https://doi.org/10.1016/j.adapen.2023.100135}

\bibitem{ref14}
Nguyen, A.-T., Reiter, S., \& Rigo, P. (2014). A review on simulation-based optimization methods applied to building performance analysis. \textit{Applied Energy}, 113, 1043--1058. \url{https://doi.org/10.1016/j.apenergy.2013.08.061}

\bibitem{ref15}
Machairas, V., Tsangrassoulis, A., \& Axarli, K. (2014). Algorithms for optimization of building design: A review. \textit{Renewable and Sustainable Energy Reviews}, 31, 101--112. \url{https://doi.org/10.1016/j.rser.2013.11.036}

\bibitem{ref19}
Zhao, Z., Li, H., \& Wang, S. (2025). Surrogate-assisted coordinated design optimization of building and microclimate considering their mutual impacts. \textit{Applied Energy}, 383. \url{https://doi.org/10.1016/j.apenergy.2025.125374}

\bibitem{ref20}
Westermann, P., \& Evins, R. (2019). Surrogate modelling for sustainable building design: A review. \textit{Energy and Buildings}, 198, 170--186. \url{https://doi.org/10.1016/j.enbuild.2019.05.057}

\bibitem{wang2026surrogateDistrictMorphology}
Wang, H., Wu, Z., Gu, W., Liu, P., Sun, Q., \& Wang, W. (2026). A surrogate-assisted framework for district-scale urban morphology optimization toward reduced building energy demand. \textit{Applied Energy}, 422, 128294. \url{https://doi.org/10.1016/j.apenergy.2026.128294}

\bibitem{ref9}
Eltaweel, A., \& Su, Y. (2017). Parametric design and daylighting: A literature review. \textit{Renewable and Sustainable Energy Reviews}, 73, 1086--1103. \url{https://doi.org/10.1016/j.rser.2017.02.011}

\bibitem{ref10}
Longo, S., Montana, F., \& Riva Sanseverino, E. (2019). A review on optimization and cost-optimal methodologies in low-energy buildings design and environmental considerations. \textit{Sustainable Cities and Society}, 45, 87--104. \url{https://doi.org/10.1016/j.scs.2018.11.027}

\bibitem{ref11}
Ascione, F., De Masi, R. F., et al. (2016). Optimization of building envelope design for nZEBs in Mediterranean climate: Performance analysis of a residential case study. \textit{Applied Energy}, 183, 938--957. \url{https://doi.org/10.1016/j.apenergy.2016.09.027}

\bibitem{ref12}
Li, G., He, Q., et al. (2025). Accelerated inverse urban design: A multi-objective optimization method to photovoltaic power generation potential, environmental performance and economic performance in urban blocks. \textit{Sustainable Cities and Society}, 120. \url{https://doi.org/10.1016/j.scs.2025.106135}

\bibitem{ref16}
Zhang, Y., Teoh, B. K., \& Zhang, L. (2024). Multi-objective optimization for energy-efficient building design considering urban heat island effects. \textit{Applied Energy}, 376. \url{https://doi.org/10.1016/j.apenergy.2024.124117}

\bibitem{ref17}
Shi, G., Yao, S., et al. (2024). Multi-performance collaborative optimization of existing residential building retrofitting in extremely arid and hot climate zones: A case study in Turpan, China. \textit{Journal of Building Engineering}, 89. \url{https://doi.org/10.1016/j.jobe.2024.109304}

\bibitem{ref23}
Liu, R., Fang, T., et al. (2024). Controllable cross-building multi-objective optimisation for NZEBs: A framework utilising parametric generation and intelligent algorithms. \textit{Applied Energy}, 374. \url{https://doi.org/10.1016/j.apenergy.2024.124003}

\bibitem{ref24}
Miraba, S., Salehi, A., et al. (2025). Adaptive BIPV shading optimization for electricity generation and building electricity management, sDA, and DGP using multi-objective algorithms. \textit{Applied Energy}, 402. \url{https://doi.org/10.1016/j.apenergy.2025.126860}

\bibitem{ref13}
Liu, K., Xu, X., et al. (2023). A multi-objective optimization framework for designing urban block forms considering daylight, energy consumption, and photovoltaic energy potential. \textit{Building and Environment}, 242. \url{https://doi.org/10.1016/j.buildenv.2023.110585}

\bibitem{ref22}
Khoroshiltseva, M., Slanzi, D., \& Poli, I. (2016). A Pareto-based multi-objective optimization algorithm to design energy-efficient shading devices. \textit{Applied Energy}, 184, 1400--1410. \url{https://doi.org/10.1016/j.apenergy.2016.05.015}

\bibitem{ref25}
Ciardiello, A., Rosso, F., et al. (2020). Multi-objective approach to the optimization of shape and envelope in building energy design. \textit{Applied Energy}, 280. \url{https://doi.org/10.1016/j.apenergy.2020.115984}

\bibitem{ref26}
Wang, X., Wang, S., et al. (2024). Deep reinforcement learning: A survey. \textit{IEEE Transactions on Neural Networks and Learning Systems}, 35(4), 5064--5078. \url{https://doi.org/10.1109/tnnls.2022.3207346}

\bibitem{puterman1994mdp}
Puterman, M. L. (1994). \textit{Markov Decision Processes: Discrete Stochastic Dynamic Programming}. Wiley.

\bibitem{sutton2018reinforcement}
Sutton, R. S., \& Barto, A. G. (2018). \textit{Reinforcement Learning: An Introduction} (2nd ed.). MIT Press.

\bibitem{ref31}
Mnih, V., Kavukcuoglu, K., et al. (2015). Human-level control through deep reinforcement learning. \textit{Nature}, 518(7540), 529--533. \url{https://doi.org/10.1038/nature14236}

\bibitem{ref27}
Jemin, H., Joonho, L., et al. (2019). Learning agile and dynamic motor skills for legged robots. \textit{Science Robotics}, 4(26). \url{https://doi.org/10.1126/scirobotics.aau5872}

\bibitem{ref28}
Vinyals, O., Babuschkin, I., et al. (2019). Grandmaster level in StarCraft II using multi-agent reinforcement learning. \textit{Nature}, 575(7782), 350--354. \url{https://doi.org/10.1038/s41586-019-1724-z}

\bibitem{ref29}
Wang, X., Chen, W., et al. (2018). Video captioning via hierarchical reinforcement learning. In \textit{2018 IEEE/CVF Conference on Computer Vision and Pattern Recognition}, 4213--4222. \url{https://doi.org/10.1109/cvpr.2018.00443}

\bibitem{ref30}
Li, J., Yao, L., et al. (2020). Deep reinforcement learning for pedestrian collision avoidance and human-machine cooperative driving. \textit{Information Sciences}, 532, 110--124. \url{https://doi.org/10.1016/j.ins.2020.03.105}

\bibitem{ref32}
Lillicrap, T. P., Hunt, J. J., et al. (2015). Continuous control with deep reinforcement learning. \textit{arXiv preprint arXiv:1509.02971}. \url{https://doi.org/10.48550/arXiv.1509.02971}

\bibitem{ref33}
Du, Y., Zandi, H., et al. (2021). Intelligent multi-zone residential HVAC control strategy based on deep reinforcement learning. \textit{Applied Energy}, 281. \url{https://doi.org/10.1016/j.apenergy.2020.116117}

\bibitem{ref34}
Li, S., Su, S., \& Lin, X. (2025). Optimizing the hyper-parameters of deep reinforcement learning for building control. \textit{Building Simulation}, 18(4), 765--789. \url{https://doi.org/10.1007/s12273-025-1233-y}

\bibitem{ref35}
Wu, Z., Li, Y., et al. (2025). Distributed voltage control for multi-feeder distribution networks considering transmission network voltage fluctuation based on robust deep reinforcement learning. \textit{Applied Energy}, 379. \url{https://doi.org/10.1016/j.apenergy.2024.124984}

\bibitem{ref36}
Kou, P., Liang, D., et al. (2020). Safe deep reinforcement learning-based constrained optimal control scheme for active distribution networks. \textit{Applied Energy}, 264. \url{https://doi.org/10.1016/j.apenergy.2020.114772}

\bibitem{ref37}
Biemann, M., Scheller, F., Liu, X., \& Huang, L. (2021). Experimental evaluation of model-free reinforcement learning algorithms for continuous HVAC control. \textit{Applied Energy}, 298, 117164. \url{https://doi.org/10.1016/j.apenergy.2021.117164}

\bibitem{ref38}
Alabi, T. M., Lawrence, N. P., Lu, L., Yang, Z., \& Gopaluni, R. B. (2023). Automated deep reinforcement learning for real-time scheduling strategy of multi-energy system integrated with post-carbon and direct-air carbon captured system. \textit{Applied Energy}, 333, 120633. \url{https://doi.org/10.1016/j.apenergy.2022.120633}

\bibitem{ref39}
Deng, X., Zhang, Y., Jiang, Y., Zhang, Y., \& Qi, H. (2024). A novel operation method for renewable building by combining distributed DC energy system and deep reinforcement learning. \textit{Applied Energy}, 353(B), 122188. \url{https://doi.org/10.1016/j.apenergy.2023.122188}

\bibitem{ref40}
Aubert, H., \& Bressel, M. (2025). Optimal sizing of isolated photovoltaic-hydrogen microgrids using covariance matrix adaptation evolution strategy considering real-gas modeling of hydrogen. \textit{Applied Energy}, 401(B), 126677. \url{https://doi.org/10.1016/j.apenergy.2025.126677}

\end{thebibliography}
